\section{Chiral currents and Gaussian correlations}\label{sec:chiral}

We now turn the angular spectrum into limits of height derivatives. The
argument has two steps. Separating two angular scales gives two independent
chiral current forms. A second separation, with one scale tending to zero,
determines the current collision coefficient. Holomorphy in both axial
quantizations then determines every separated current correlation from this
coefficient.

We work at a fixed parameter in the new range and use the diagonal-word
Hilbert space $\mathcal D$ of Section~\ref{sec:angular}. Its vacuum is
$\Omega$. The commuting rows $P(w)$ have spectral functions $f$, the marked
row is $J(w)$, and
\[
 H_0=-\log|P(1)|,\qquad S=P(i),\qquad
 N(w)=\frac{J(w)P(w)+P(w)J(w)}2.
\]
On this space, Lemma~\ref{lem:inner-spectra} gives finite Blaschke products
with $f(0)=f(\infty)=1$ and real symmetry, including the constant product
$f=1$. Their zeros $\zeta$ lie in the fixed cone
$\Re\zeta\ge c_0|\zeta|$. The current vacuum identity and its measure are
\begin{equation}\label{eq:chiral-vacuum-data}
 J(w)\Omega=(1-f(w))\mathsf h,\qquad
 \mu(df)=\|\mathsf h(f)\|^2\sigma(df),\qquad
 \mu=ds\,\nu,\qquad \nu(1)=\kappa<\infty.
\end{equation}
Here $\sigma$ is a scalar spectral measure and $s$ is logarithmic
dilation along a nonconstant spectral orbit.
The symbol $\mathsf h$ is a spectral field; it need not be a vector of finite
norm. Every use of it below includes an integrable multiplier.

\subsection{The two chiral forms}

Define nonnegative spectral multiplication operators on $\mathcal D$ by
\begin{equation}\label{eq:chiral-energies}
 a_+(f)=-f'(0)=\sum_{\zeta}2\Re(\zeta^{-1}),\qquad
 a_-(f)=-\left.\frac{d}{dv}f(v^{-1})\right|_{v=0}
                         =\sum_{\zeta}2\Re\zeta.
\end{equation}
Both values are zero on $f=1$, and both are strictly positive on every
nonconstant fiber. Write $E_\pm(R)=\bbone_{a_\pm\le R}$ and
$\mathcal D_\pm^{\mathrm{fin}}=\bigcup_{R>0}E_\pm(R)\mathcal D$.
These are dense linear subspaces. A current form on this domain means a
sesquilinear form whose restriction to each bounded spectral band is a
bounded operator; no unbounded-operator closure is implicit in this usage.

\begin{lemma}\label{lem:chiral-forms}
There are consistent skew-adjoint forms $j_\pm$ on
$\mathcal D_\pm^{\mathrm{fin}}$ such that, for every $R>0$,
\begin{equation}\label{eq:chiral-form-bound}
 \|E_\pm(R)j_\pm E_\pm(R)\|\le CR,
 \qquad E_\pm(R)j_\pm\Omega=E_\pm(R)a_\pm\mathsf h.
\end{equation}
Their vacuum means vanish. For $s>0$, the damped forms extend to bounded
operators and satisfy
\begin{equation}\label{eq:chiral-damped-form}
                     \|e^{-sa_\pm}j_\pm e^{-sa_\pm}\|\le C/s.
\end{equation}
For every nonnegative measurable function $g$ on $(0,\infty)$,
\begin{equation}\label{eq:chiral-Haar-integral}
             \int g(a_\pm(f))\dd\mu(f)
                     =\kappa\int_0^\infty g(t)\frac{dt}{t}.
\end{equation}
\end{lemma}

\begin{proof}
We give the construction at zero. The construction at infinity uses the
coordinate $v=1/w$ and the same estimates. The reciprocal-radius budget
in \eqref{eq:chiral-energies}, together with the cone condition, gives
uniformly on the band $E_+(R)$
\begin{equation}\label{eq:chiral-disc-inverse}
 P(w)=I+O(R|w|),\qquad \|P(w)^{-1}\|\le C,
                  \quad \Re w\ge0,\quad |w|<c/R.
\end{equation}
For example, normalize each Blaschke factor to one at zero. Its difference
from one is at most $C|w|\Re(\zeta^{-1})$ in this disc, and the sum of
these bounds is at most $CR|w|$. Taking $c$ small bounds the product and
its inverse. All inverses in this proof are restricted to the stated band.

On $E_+(R)\mathcal D$ write $P_R(w)$ for the restriction of $P(w)$.
Compress $J$ by $E_+(R)$ and continue it to the other half-disc by the
reversal identity
\begin{equation}\label{eq:chiral-gluing}
 J_R(-w)=-P_R(w)^{-1}J_R(w)P_R(w)^{-1},
       \qquad J_R=E_+(R)JE_+(R).
\end{equation}
For this boundary gluing, first dilate the uncompressed endpoint reversal
identity and then compress by $E_+(R)$, which commutes with every $P(w)$.
Thus the argument of Lemma~\ref{lem:low-band} applies without moving a
cutoff through a boost. The resulting
function is bounded and holomorphic on the punctured disc. Its center is
removable: the negative Laurent coefficients vanish by the Cauchy bounds.
At the center \eqref{eq:chiral-gluing} says $J_R(0)=-J_R(0)$, so that
constant coefficient is zero. Define $E_+(R)j_+E_+(R)$ to be the linear
Taylor coefficient. The Cauchy estimate on a circle of radius proportional
to $R^{-1}$ proves the bound $CR$.

Compression from a larger band to a smaller band commutes with taking
this coefficient. This proves consistency and defines the form on the
union of bands. For positive real $w$, $J(w)^*=-J(w)$, so the coefficient
is skew-adjoint as a form. The same argument gives $j_-$.

Dilation multiplies $a_+$ by the dilation factor and $a_-$ by its
reciprocal. On each nonconstant orbit either substitution changes Haar
measure $ds$ to $dt/t$. Equation~\eqref{eq:chiral-vacuum-data} therefore
gives \eqref{eq:chiral-Haar-integral}. In particular,
$\int_{a_\pm\le R}a_\pm^2\dd\mu=\kappa R^2/2$.
Dividing the vacuum identity by $w$, or by $v$ at infinity, and taking
the band Taylor coefficient now gives the vacuum assertion in
\eqref{eq:chiral-form-bound}. Its inner product with $\Omega$ is zero:
$\Omega$ lies in the zero-energy band, and its compressed current mean
is bounded by $CR$ for every $R>0$.

For the damping assertion, divide the spectrum of $a_\pm$ into
$[0,s^{-1}]$ and $(2^{q-1}s^{-1},2^qs^{-1}]$, $q\ge1$.
The form bound controls a block between levels $q,r$ by
$Cs^{-1}\max(2^q,2^r)$. Damping on both sides supplies
$e^{-2^{q-1}}e^{-2^{r-1}}$ when the respective index is positive.
The sum of these block norms is at most $C/s$, which proves both
existence of the bounded extension and \eqref{eq:chiral-damped-form}.
The same argument shows that tails of these damped block sums tend to
zero in operator norm.
\end{proof}

\subsection{Separated scales and representative norms}

Let $\mathcal C$ be the span of strict diagonal words applied to $\Omega$.
For two positive scales $r_1,r_2$, represent an elementary tensor
$u_1\otimes u_2$ of such words in $\mathcal D$ by replacing every argument
in word $i$ by $r_iw$, multiplying the two resulting words in that order,
and applying the product to $\Omega$. Extend this prescription to finite
linear combinations and denote a chosen representative of $u$ by
$u_{r_1,r_2}$. Different word expressions for the same tensor have
difference tending to zero in norm as $|\log(r_1/r_2)|\to\infty$.
Indeed, Proposition~\ref{prop:word-clustering}, applied after adjunction,
gives
\begin{equation}\label{eq:chiral-Gram-limit}
       \langle u_{r_1,r_2},v_{r_1,r_2}\rangle
                          \longrightarrow\langle u,v\rangle
                  \quad(u,v\in\mathcal C\otimes\mathcal C).
\end{equation}
The asymptotic diagonal RTT exchange permits the two separated word
groups to be interchanged with an error tending to zero in norm.
Consequently a fixed bounded diagonal marked row at $r_iw$ acts in the
limit on factor $i$ alone, for every strict $w$.

We use convergence of representatives, not a bounded map from the entire
tensor completion. Precisely, a family $v_\delta\in\mathcal D$ converges
strongly to $v\in\mathcal D\otimes\mathcal D$ if, for every $\epsilon>0$,
there is a finite word tensor $u$ with
\[
 \|u-v\|<\epsilon,\qquad
 \limsup_{\delta\downarrow0}\|v_\delta-u_\delta\|<\epsilon.
\]
The scales defining $u_\delta$ will always be specified. Equation
\eqref{eq:chiral-Gram-limit} makes this definition independent of the
chosen word expressions.

\begin{lemma}\label{lem:chiral-representatives}
Suppose $T_\delta$ are uniformly bounded operators and, on every finite
word tensor, $T_\delta u_\delta$ converges strongly to $Tu$. Then $T$
extends to a bounded operator on the tensor completion, and
$T_\delta v_\delta\to Tv$ for every strongly converging family
$v_\delta\to v$. Such operator limits may be composed, without a
commutation assumption.

A family $v_\delta$ converges strongly to $v$ if its pairings with all
core representatives converge to the pairings with $v$ and its norms
converge to $\|v\|$.
\end{lemma}

\begin{proof}
If $\|T_\delta\|\le M$, then \eqref{eq:chiral-Gram-limit} implies
$\|Tu\|\le M\|u\|$ on the core. Approximate $v_\delta$ by $u_\delta$
as in the definition. The norm of
$T_\delta(v_\delta-u_\delta)$ is at most $M$ times that approximation
error, while $T_\delta u_\delta$ has the asserted limit. Taking the
approximation error to zero proves the first assertion and then composition.
For the last assertion, use
\[
 \|v_\delta-u_\delta\|^2
 =\|v_\delta\|^2+\|u_\delta\|^2
       -2\Re\langle v_\delta,u_\delta\rangle
                         \longrightarrow\|v-u\|^2
\]
and approximate $v$ by a core vector $u$.
\end{proof}

\begin{proposition}[Two-scale spectral limit]\label{prop:two-scale}
Take $r_1=\delta^{-1}$ and $r_2=\delta$, with $\delta\downarrow0$ through
arbitrary positive real values. On core representatives, the joint spectral
functions $(f(\delta^{-1}\,\cdot),f(\delta\,\cdot))$ converge in law in
the compact-open topology of the right half-plane to the two tensor-factor
functions $(F_1,F_2)$. Jointly with these functions,
\begin{equation}\label{eq:chiral-scaled-spectrum}
\begin{split}
 \delta^{-1}H_0&\longrightarrow A:=a_+\otimes I+I\otimes a_-,\\
 \delta^{-1}\arg S&\longrightarrow -a_+\otimes I+I\otimes a_-,
       \qquad \sgn P(1)\longrightarrow I
\end{split}
\end{equation}
in spectral law. The argument is the principal argument near zero.
There is an unbounded set of $R>0$ for which
\begin{equation}\label{eq:chiral-band-limit}
 Q_\delta(R):=\bbone_{H_0/\delta\le R}
                       \longrightarrow Q(R):=\bbone_{A\le R}
\end{equation}
strongly on representatives. Bounded spectral multipliers obtained jointly
from these limits have the same strong interpretation whenever their
discontinuity sets have zero limiting state mass.

If $m_\delta\ge0$ and $l_\delta$ are integers with
$\delta m_\delta\to x\ge0$ and $\delta l_\delta\to y$, then
\begin{equation}\label{eq:chiral-semigroup-limit}
 P(1)^{m_\delta}S^{l_\delta}\longrightarrow
       e^{-(x+iy)a_+}\otimes e^{-(x-iy)a_-}
\end{equation}
strongly on representatives.
\end{proposition}

\begin{proof}
The strong row actions preceding Lemma~\ref{lem:chiral-representatives}
give convergence of every polynomial in finitely many samples of the two
spectral functions and their complex conjugates. The functions are bounded
by one on the right half-plane; normal-family compactness therefore gives
the joint compact-open convergence. The limiting functions are the original
finite Blaschke spectral functions on the two factors, and are not
identically zero. For a vector $v$, its state spectral measure is
$\sigma_v(B)=\|\bbone_Bv\|^2$ on measurable sets of spectral functions.
These measures have total masses
$\|u_\delta\|^2\to\|u\|^2$; normalization to probability is harmless
when $u\ne0$.

The rescaled quantities in \eqref{eq:chiral-scaled-spectrum} require
control of \emph{every} zero, including zeros between the two scales.
For a fixed core representative, the exterior RTT estimate from
Lemma~\ref{lem:inner-spectra} gives, outside fixed word offsets,
\begin{equation}\label{eq:chiral-exterior-dip}
 \|(I-P(e^s))u_{r_1,r_2}\|^2
      \le C_u\sum_{i=1}^2e^{-2\eta|s-\log r_i|}.
\end{equation}
Every zero $\zeta$ forces a fixed-width interval about $\log|\zeta|$
on which $|1-f(e^s)|$ is bounded below by a positive constant. This
follows from the cone condition and the modulus loss of its single
Blaschke factor; all other factors have modulus at most one.
Integrating \eqref{eq:chiral-exterior-dip} outside slightly smaller
windows thus yields
\begin{equation}\label{eq:chiral-all-zero-event}
\begin{split}
 &\sigma_{u_{r_1,r_2}}
 \left\{\text{some zero lies outside }
     \bigcup_{i=1}^2\{\zeta:|\log(|\zeta|/r_i)|\le D\}\right\}\\
 &\hspace{40mm}\le C_u e^{-cD}.
\end{split}
\end{equation}
An additive fixed change in $D$ absorbs the dip width and the word offsets.
To justify the event bound explicitly, if any zero lies outside the enlarged
windows, one entire dip interval lies in the integration region. The
integral of $|1-f(e^s)|^2$ there is then at least a fixed positive number.
Markov's inequality after integration over the state spectral measure gives
\eqref{eq:chiral-all-zero-event}.

Here is the order of truncations. Given a desired exceptional mass, first
choose $D$ large enough in \eqref{eq:chiral-all-zero-event} and so that
the finite zero sets of $F_1,F_2$ lie in their annuli with correspondingly
high probability. Choose their boundary circles to carry no limiting zero
almost surely. Only countably many radii are excluded: first restrict to
a bounded zero count, then apply this assertion to its finite intensity
measure, and finally take the union over count bounds. On these fixed
annuli, the cone keeps the zeros in a compact subset of the open right
half-plane. The argument principle and compact-open convergence give joint
convergence of the complete zero multisets, including multiplicities.
Next choose a finite degree bound $M$ so that larger limiting counts have
arbitrarily small mass. Take the scale limit only after $D$ and $M$ have
been fixed. Zero counts and locations now lie in a compact, bounded-count
set, so every product expansion below is uniform. Finally let the allowed
exceptional mass tend to zero, increasing $D$ and $M$ if necessary.

Each limiting product has even degree: conjugate zeros pair, real zeros
contribute a factor $-1$ between the two ends, and both end values are one.
On the preceding high-mass event the complete limiting zero counts occur
in the chosen annuli. Multiplicity convergence therefore makes both
finite-scale clusters even with probability tending to one. This proves
the needed parity; it is not an assumption on a truncated cluster.

For clarity, write the large zeros as $\delta^{-1}\alpha_j$ and the
small zeros as $\delta\beta_j$. On the retained event each normalized
zero has modulus between $e^{-D}$ and $e^D$, and there are at most $M$
zeros. Normalize each factor to one at zero. The two even clusters then
give, for $w=1,i$,
\begin{equation}\label{eq:chiral-product-expansion}
 f(w)=1-\delta\left(
       w\sum_j2\Re\alpha_j^{-1}
          +w^{-1}\sum_j2\Re\beta_j\right)+O_{D,M}(\delta^2).
\end{equation}
The two sums converge jointly to $a_+(F_1)$ and $a_-(F_2)$.
Taking logarithmic modulus at $1$ and phase at $i$ gives
\eqref{eq:chiral-scaled-spectrum}. The middle sign is positive on this
event. Values of $H_0/\delta$ can be arbitrarily large on the discarded
event; convergence in law and bounded multiplier convergence require no
first-moment estimate for that quantity.

We next strengthen bounded spectral convergence to convergence of vectors.
Choose a countable total word core, including its rational linear
combinations, and exclude the union of the atoms of $A$ in all their state
measures. That union is countable; any other $R$ is admissible in
\eqref{eq:chiral-band-limit}. At such $R$, continuous cutoffs approximate
the indicator in each state $L^2$ norm. More generally, let $b_\delta$ be
a uniformly bounded spectral multiplier with the stated joint limit $b$.
The algebra of polynomials in a countable dense set of interior samples
and their conjugates generates the spectral functions and is dense in
the state $L^2$ space. Given a state $u$, choose such a polynomial $p$ with
$\|(b-p)u\|$ small. Joint convergence gives
\begin{equation}\label{eq:chiral-multiplier-norm}
 \|(b_\delta-p_\delta)u_\delta\|^2
          \longrightarrow\|(b-p)u\|^2.
\end{equation}
For indicators, take the continuous approximations first. Polynomial
actions already converge strongly. Equation
\eqref{eq:chiral-multiplier-norm} proves strong convergence of the
multiplier actions; polarization gives the corresponding mixed pairings.
Lemma~\ref{lem:chiral-representatives} extends these actions to all
strongly converging families and permits their compositions. No assertion
that a current preserves the exceptional-zero event is needed.

Finally, $|f(1)|\le1$ and $|f(i)|=1$. The bounded multipliers
$f(1)^{m_\delta}f(i)^{l_\delta}$ converge jointly by
\eqref{eq:chiral-scaled-spectrum} to
$e^{-xA+iy(-a_+\otimes I+I\otimes a_-)}$, which is
\eqref{eq:chiral-semigroup-limit}. The same vector argument applies.
\end{proof}

\begin{figure}[htbp]
\centering
\input{figures/chiral-windows}
\caption{The two windows in logarithmic radius. Their width and a bound
on the number of zeros are fixed before $\delta$ tends to zero. Outside
an event of arbitrarily small state mass, no zero lies between or beyond
the windows.}
\label{fig:chiral-windows}
\end{figure}

The exceptional-event estimate is the reason that local spectral convergence
can be used at the intervening physical scale $w=1$. It also permits the
energy cutoffs needed for current insertions. We next prove their convergence
in norm, before removing those cutoffs by time damping.

\subsection{Current insertions and damping}

\begin{lemma}\label{lem:current-limit}
For every admissible $R>0$ in Proposition~\ref{prop:two-scale} and every
$|w|=1$ with $\Re w\ge0$,
\begin{equation}\label{eq:chiral-current-limit}
 \delta^{-1}Q_\delta(R)N(w)Q_\delta(R)
 \longrightarrow
 Q(R)\bigl(w(j_+\otimes I)+w^{-1}(I\otimes j_-)\bigr)Q(R)
\end{equation}
strongly on representatives. All displayed operators are bounded at fixed
$R$. Finite products with propagators from
\eqref{eq:chiral-semigroup-limit} converge when successive positive time
gaps stay bounded below. In these vacuum products, the cutoffs on both
sides of each current may be removed after the scale limit, with error
tending to zero uniformly on compact sets of separated time coordinates.
\end{lemma}

\begin{proof}
Write $Q_\delta=Q_\delta(R)$ and $Q=Q(R)$. Glue
$J_\delta(w)=Q_\delta J(w)Q_\delta$ to the full annulus
\[
                  CR\delta<|w|<c/(R\delta)
\]
using Lemma~\ref{lem:low-band}. Let $C_{k,\delta}$ be its Laurent
coefficients. Fix $\rho>0$ sufficiently small compared with $R^{-1}$.
At the outer circle, the substitution $v=\delta w$ gives
\begin{equation}\label{eq:chiral-outer-coefficient}
 \delta^{-1}C_{1,\delta}
    =\frac1{2\pi i}\int_{|v|=\rho}
                  J_\delta(v/\delta)\frac{dv}{v^2}.
\end{equation}
At the inner circle, put $w=\delta v$ to obtain
\begin{equation}\label{eq:chiral-inner-coefficient}
 \delta^{-1}C_{-1,\delta}
    =\frac1{2\pi i}\int_{|v|=\rho^{-1}}
                              J_\delta(\delta v)\,dv.
\end{equation}
Both contours are counterclockwise. The rescaled contours are fixed, and
both integrands are uniformly bounded in operator norm at fixed $R,\rho$.

On the strict right semicircle the marked rows converge on their respective
tensor factors. The cutoffs converge strongly by
\eqref{eq:chiral-band-limit}. On the reflected semicircle the integrand is
the product specified by reversal, for example
\[
 -\bigl(Q_\delta P(v/\delta)^{-1}Q_\delta\bigr)
     J_\delta(v/\delta)
     \bigl(Q_\delta P(v/\delta)^{-1}Q_\delta\bigr).
\]
The inverse factors are uniformly bounded spectral multipliers on
$Q_\delta\mathcal D$ only. Their joint limits follow from
Proposition~\ref{prop:two-scale}, since on $Q\mathcal D^{\otimes2}$
the first energy is at most $R$ and \eqref{eq:chiral-disc-inverse}
applies. The second contour uses the analogous estimate at infinity.
Lemma~\ref{lem:chiral-representatives} therefore proves convergence of
these noncommuting products at every strict contour point.

Here is the norm justification for passing through the contour integrals.
For a representative $u_\delta$, let $v_\delta(t)$ be the integrand
applied to it, including the scalar contour weight, and let $v(t)$ be
its limit. The vectors are bounded uniformly in $t,\delta$. Core test
pairings of $\int v_\delta(t)dt$ converge by a single dominated integral.
Their squared norms converge because
\begin{equation}\label{eq:chiral-double-contour}
 \left\|\int v_\delta(t)dt\right\|^2
 =\iint\langle v_\delta(t),v_\delta(s)\rangle\,dt\,ds
 \longrightarrow
 \iint\langle v(t),v(s)\rangle\,dt\,ds.
\end{equation}
Pairings of strongly convergent representatives converge, so dominated
convergence applies to this double integral. The two imaginary-axis points
have contour measure zero. The norm criterion in
Lemma~\ref{lem:chiral-representatives} now gives strong convergence of
each coefficient action. By the definitions of the chiral forms, the
limits of \eqref{eq:chiral-outer-coefficient} and
\eqref{eq:chiral-inner-coefficient} are respectively
\[
                 Q(j_+\otimes I)Q,
                   \qquad Q(I\otimes j_-)Q.
\]

It remains to pass from these coefficients to $N$. Put
$\Sigma=\sgn P(1)$ and
\[
 T_\delta(w)=w\delta^{-1}C_{1,\delta}
                         +w^{-1}\delta^{-1}C_{-1,\delta}.
\]
The same-sign and opposite-sign estimates of Lemma~\ref{lem:low-band}
give, for some $\theta>0$,
\begin{equation}\label{eq:chiral-sign-removal}
 \delta^{-1}Q_\delta N(w)Q_\delta
       =\tfrac12\bigl(\Sigma T_\delta(w)+T_\delta(w)\Sigma\bigr)
                       +O_R(\delta^\theta)
\end{equation}
in operator norm, uniformly on the unit semicircle. Indeed the
anticommutator is $+T_\delta$ on the positive-sign block,
$-T_\delta$ on the negative-sign block, and zero between opposite signs.
On same-sign blocks the even Laurent part and all odd modes of order
larger than one vanish after rescaling, with the stated remainder.
The opposite-sign block of the rescaled $N$ has the same vanishing bound.
Proposition~\ref{prop:two-scale} gives $\Sigma\to I$ strongly on
representatives. The coefficient operators are uniformly bounded and have
the strong limits just proved. Composition in
\eqref{eq:chiral-sign-removal} proves \eqref{eq:chiral-current-limit}.
The low-band bounds include the endpoints $w=\pm i$.

We finish with a summable bound for removing cutoffs. For all $U>0$ and
all sufficiently small $\delta$,
\begin{equation}\label{eq:chiral-all-band-bound}
 \|\bbone_{H_0/\delta\le U}\,\delta^{-1}N(w)
                      \bbone_{H_0/\delta\le U}\|\le CU.
\end{equation}
If $U\delta$ is small this is Lemma~\ref{lem:low-band}; otherwise the
global bound on $N$ gives it by $\delta^{-1}\le C U$.
Decompose $H_0/\delta$ into bands $[0,1]$ and
$(2^{q-1},2^q]$, $q\ge1$. A current block between levels $q,r$ costs
at most $C\max(2^q,2^r)$. Each internal propagator with
$\delta m\ge b>0$ supplies at level $q\ge1$ the factor
$e^{-b2^{q-1}}$, since
$\|P(1)^m\bbone_{H_0/\delta>2^{q-1}}\|\le e^{-b2^{q-1}}$;
unit shifts do not affect this norm.
The outer states are the vacuum, of energy zero. For a product of $k$
currents, a bound for the sum over its $k-1$ internal indices is a constant
times
\begin{equation}\label{eq:chiral-dyadic-majorant}
 \sum_{q_1,\ldots,q_{k-1}\ge0}
       \prod_{j=1}^{k-1}(1+2^{q_j})^2e^{-b'2^{q_j}},
                         \qquad b'>0.
\end{equation}
This converges, and its tails with any $2^{q_j}>R$ tend to zero uniformly
in the mesh. The limiting forms have the identical bound using $A$ and
\eqref{eq:chiral-form-bound}. Fixed-cutoff convergence, followed by this
tail estimate, proves the product assertion. The constants are uniform
when the finitely many limiting time gaps have a fixed positive lower
bound.
\end{proof}

\subsection{Continuum derivative correlations}

In the remainder of this section $z=x+iy$ denotes a physical point, not
a row angle. At a mesh face define the centered height differences
\begin{align*}
 D_{x,\delta}h_\delta(z)
   &=\tfrac12\bigl(h_\delta(z+\delta)-h_\delta(z-\delta)\bigr),\\
 D_{y,\delta}h_\delta(z)
   &=\tfrac12\bigl(h_\delta(z+i\delta)-h_\delta(z-i\delta)\bigr),\\
 D_\delta^+&=\tfrac12(D_{x,\delta}-iD_{y,\delta}),\qquad
 D_\delta^-=\tfrac12(D_{x,\delta}+iD_{y,\delta}).
\end{align*}
The mesh inverse $\delta^{-1}$ is not included in these definitions.
One may evaluate at any point using the stipulated face convention.
All arguments below allow the chosen lattice faces to move with the mesh.

Take points whose real coordinates are pairwise distinct. In eastward
quantization, a centered horizontal or vertical difference is the mean of
two consecutive height steps. The path convention and $N=(JP+PJ)/2$
identify its insertion as $-N(1)$ or $-N(i)$, respectively. To read a
product, join these two-step segments in increasing time order, using
positive time steps and signed transverse steps. Between their endpoints
the propagators are $P(1)^{m_\delta}S^{l_\delta}$. The endpoint offsets
change $m_\delta,l_\delta$ by at most a fixed number of sites and do not
change their rescaled limits.

Lemma~\ref{lem:current-limit} and
\eqref{eq:chiral-semigroup-limit} show that the insertions of
$\delta^{-1}D_\delta^+$ and $\delta^{-1}D_\delta^-$ converge to
$-j_+\otimes I$ and $-I\otimes j_-$, respectively. For example,
the two combinations follow from
\[
 N(1)\rightsquigarrow j_+\otimes I+I\otimes j_-,
 \qquad
 N(i)\rightsquigarrow i j_+\otimes I-i I\otimes j_-.
\]
Currents in different factors commute. Their propagators and the tensor
vacuum factor as well. The limiting mixed correlations therefore factor
into a pure $+$ correlation and a pure $-$ correlation.

Write $a=a_+$ and $j=j_+$ for the moment. For increasing real parts,
$\Re z_{i_1}<\cdots<\Re z_{i_k}$, define
\begin{equation}\label{eq:chiral-ordered-function}
 F_k(z_1,\ldots,z_k)
 =\left\langle\Omega,\,
 (-j)e^{-(z_{i_k}-z_{i_{k-1}})a}(-j)\cdots
 e^{-(z_{i_2}-z_{i_1})a}(-j)\Omega\right\rangle.
\end{equation}
Set $F_0=1$ and $F_1=0$. The products in
\eqref{eq:chiral-ordered-function} are defined by inserting positive
real damping on both sides of each form; arbitrary extra end damping
may be inserted because $a\Omega=0$. Lemma~\ref{lem:chiral-forms}
makes the resulting factors bounded, and the dyadic tail bound shows
that the value is independent of this allocation. By reality of the
height field, the pure $-$ answer is $\overline{F_k(z_1,\ldots,z_k)}$.
All the preceding convergence is uniform on compact subsets of the
region of distinct real coordinates. Indeed, it holds along every
sequence of moving mesh faces converging to such a configuration,
and the cutoff and damping bounds are uniform on its compact subsets.

\begin{lemma}\label{lem:chiral-holomorphy}
Let
\[
 U_x=\{(z_1,\ldots,z_k):\Re z_i\ne\Re z_j\text{ for }i\ne j\},
 \quad
 U_y=\{(z_1,\ldots,z_k):\Im z_i\ne\Im z_j\text{ for }i\ne j\}.
\]
The function $F_k$ has a symmetric, jointly holomorphic extension to
$U_x\cup U_y$, and is the limit of the centered $+$ derivative
correlations on both sets.

Fix $z_2,\ldots,z_k$ with distinct real coordinates and distinct imaginary
coordinates. As a function of $z_1$, the extension is holomorphic away
from the grid points
\begin{equation}\label{eq:chiral-grid-points}
                    p_{jh}=\Re z_j+i\Im z_h,\qquad 2\le j,h\le k.
\end{equation}
Near any such point it is $O(|z_1-p_{jh}|^{-2})$, locally uniformly
as the fixed points vary while retaining these coordinate separations.
For $k\ge2$ it tends to zero as $|z_1|\to\infty$ with the other points fixed.
\end{lemma}

\begin{proof}
On a compact subset of an eastward ordering chamber, give each current
a fixed positive amount of damping on both sides. The remaining
semigroups have parameters with real parts bounded below. These
semigroups are holomorphic in operator norm there: for instance
$\|a^m e^{-sa}\|\le m!s^{-m}$ after reducing $s$ by a fixed factor.
The resulting bounded product is jointly holomorphic in all its
parameters. This proves joint holomorphy on $U_x$, rather than merely
separate holomorphy.

By Proposition~\ref{prop:physical-plane}, northward quantization describes
the same physical law. In its coordinates $z'=-iz$, the centered
differences transform as
$D_{x',\delta}=D_{y,\delta}$ and
$D_{y',\delta}=-D_{x,\delta}$, hence
$D_\delta^+=-iD_\delta^{\prime+}$. The possible face-origin shift is
one lattice spacing and is covered by moving-face convergence. The
northward formula is thus holomorphic on $U_y$. On $U_x\cap U_y$
both formulas are limits of the same physical moments, so they agree
and glue. Permuting the insertions does not change a physical product,
which proves symmetry.

With the other points as in the statement, failure of both axial
orderings occurs only at the grid points \eqref{eq:chiral-grid-points}.
Choose a small neighborhood of the fixed configuration in which all
its real gaps and imaginary gaps remain at least $b>0$. Near $p_{jh}$,
the eastward representation has only one possibly small gap,
$g_x=|\Re(z_1-p_{jh})|$, between $z_1$ and $z_j$.
The northward representation has only one possibly small gap,
$g_y=|\Im(z_1-p_{jh})|$, between $z_1$ and $z_h$.
Use the ordering with $g=\max(g_x,g_y)\ge|z_1-p_{jh}|/\sqrt2$.
Only the two currents bordering that gap need shrinking damping.
Give each damping $g/4$ on both sides: the inner dampings consume
$g/2$, and the outer dampings come from neighboring fixed gaps or
from a vacuum endpoint. All other currents receive fixed damping.
Equation~\eqref{eq:chiral-damped-form} gives a cost $C/g$ for each
of the two currents and a bounded cost for the rest. This proves
the uniform quadratic bound. A point sharing the other coordinate
is separated in the chosen ordering.

Finally, when $|z_1|$ is large, one coordinate makes $z_1$ extreme in
the corresponding ordering, with a gap $g$ of order $|z_1|$ to all
other points. Give the isolated current damping of order $g$ on both
sides, using free damping on its vacuum side. Keep fixed damping
around all remaining currents. The isolated factor costs $C/g$;
imaginary semigroup factors have norm one. Thus $F_k\to0$ at infinity.
\end{proof}

\subsection{The collision operator}

The remaining input for Gaussianity is the coefficient at a genuine
collision. A second separation of scales shows that a short pair of
chiral currents becomes a scalar, even between nonvacuum exterior states.

\begin{lemma}\label{lem:collision}
For $a=a_+$ and $j=j_+$, the bounded operators
\begin{equation}\label{eq:chiral-collision-operator}
 K_\epsilon=e^{-\epsilon a}(\epsilon j)e^{-\epsilon a}
                      (\epsilon j)e^{-\epsilon a}
\end{equation}
are uniformly bounded for $\epsilon>0$ and satisfy
\begin{equation}\label{eq:chiral-collision-limit}
                    K_\epsilon\rightharpoonup-\kappa I
                       \qquad(\epsilon\downarrow0).
\end{equation}
Consequently, for generic distinct exterior points,
\begin{equation}\label{eq:chiral-double-pole}
 \lim_{\epsilon\downarrow0}\epsilon^2
 F_k(z_j+\epsilon,z_2,\ldots,z_k)
   =-\kappa F_{k-2}(z_2,\ldots,\widehat z_j,\ldots,z_k).
\end{equation}
The coefficient on the right is independent of both colliding positions.
\end{lemma}

\begin{proof}
Split the middle damping of \eqref{eq:chiral-collision-operator} into
two halves. Each of the resulting damped current factors has norm bounded
by a constant, by \eqref{eq:chiral-damped-form} at scale $\epsilon/2$.
This gives uniform boundedness.

Apply the representative construction with scales $r_1=\epsilon$ and
$r_2=1$. Its Gram and row limits use only separation of the two scales,
so they are unchanged. The all-zero event estimate
\eqref{eq:chiral-all-zero-event}, with the same order of fixed windows,
bounded degree, and then scale limit, gives
\begin{equation}\label{eq:chiral-collision-energy}
                     \epsilon a(f)\longrightarrow a(F_1)
\end{equation}
jointly with the two spectral functions. To see the scaling directly,
a zero $\epsilon\alpha$ contributes $2\Re\alpha^{-1}$ to $\epsilon a$,
whereas a zero $\beta$ at the unshifted scale contributes
$2\epsilon\Re\beta^{-1}$. On the bounded-count annuli the latter
sum tends uniformly to zero. The exceptional-zero event has arbitrarily
small mass and is discarded before making this expansion.

Choose a non-atomic band endpoint and put
$B_\epsilon=\bbone_{\epsilon a\le R}$ and
$B=\bbone_{a\otimes I\le R}$. The spectral multiplier argument proves
$B_\epsilon\to B$ strongly. The definition of $j$ by its Taylor
coefficient gives, for a sufficiently small fixed $\rho<c/R$,
\begin{equation}\label{eq:chiral-collision-contour}
 \epsilon B_\epsilon jB_\epsilon
   =\frac1{2\pi i}\int_{|v|=\rho}
             B_\epsilon J(\epsilon v)B_\epsilon\,\frac{dv}{v^2}
   \longrightarrow B(j\otimes I)B.
\end{equation}
Indeed, the unscaled band is $a\le R/\epsilon$, whose Taylor disc
has radius proportional to $\epsilon/R$. On the right semicircle
the row acts on the first factor; on the left semicircle the bounded
compressed inverse and reversal apply. The single- and double-contour
argument \eqref{eq:chiral-double-contour} proves convergence in
representative norm. Thus \eqref{eq:chiral-collision-energy} and
\eqref{eq:chiral-collision-contour} give the cutoff limit of
$K_\epsilon$.

These cutoffs can be removed in operator norm uniformly in $\epsilon$.
Decompose $\epsilon a$ into $[0,1]$ and the dyadic bands above one.
A block of $\epsilon j$ between levels $q,r$ costs at most
$C\max(2^q,2^r)$. The three damping factors of $K_\epsilon$ give an
absolutely summable majorant
\[
 C\sum_{q_0,q_1,q_2\ge0}
   \max(2^{q_0},2^{q_1})\max(2^{q_1},2^{q_2})
                       e^{-b(2^{q_0}+2^{q_1}+2^{q_2})}.
\]
Absorbing the lowest-band constants into $C$ gives some fixed $b>0$.
The part with any $2^{q_i}>R$ tends to zero as $R\to\infty$.
The same estimate holds for the limiting tensor product. Therefore,
on separated-scale representatives,
\[
 K_\epsilon\longrightarrow
           \bigl(e^{-a}je^{-a}je^{-a}\bigr)\otimes I.
\]

A fixed unboosted word vector $u$ is exactly the representative of
$\Omega\otimes u$: its small-scale word is the identity. Hence for
fixed word vectors $u,v$ the last limit contracts the first factor to
\[
 \langle\Omega,je^{-a}j\Omega\rangle
   =-\int a(f)^2e^{-a(f)}\dd\mu(f)
   =-\kappa\int_0^\infty t e^{-t}\,dt=-\kappa.
\]
The sign uses the skew-adjointness of $j$; the integrals are justified
by \eqref{eq:chiral-Haar-integral}. This proves
\eqref{eq:chiral-collision-limit} on a dense set of test vectors.
Uniform boundedness extends it to every pair of vectors in $\mathcal D$.

For \eqref{eq:chiral-double-pole}, take $z_1=z_j+\epsilon$ in the
eastward ordering. The two colliding currents are consecutive. After
multiplication by $\epsilon^2$, their part of the product becomes
$K_\epsilon$, provided its two outer dampings are subtracted from the
adjacent exterior gaps. Those gaps remain positive for small $\epsilon$;
at an endpoint, insert damping freely on the vacuum. The two exterior
vectors converge strongly, because all their damping gaps remain positive.
They may therefore be replaced by their limits in the weak convergence
of the uniformly bounded $K_\epsilon$. The two insertion minus signs
cancel. Omitting the colliding currents recombines the adjoining
semigroups, giving exactly the right side of
\eqref{eq:chiral-double-pole}. No dependence on the collision position
remains.
\end{proof}

\subsection{Removal of artificial poles and Wick recursion}

\begin{proposition}\label{prop:gaussian-currents}
Let $k\ge1$, and let $z_1,\ldots,z_k$ have pairwise distinct real
coordinates or pairwise distinct imaginary coordinates. For
$\varepsilon_i\in\{+,-\}$ and moving mesh faces $z_{i,\delta}\to z_i$,
\begin{equation}\label{eq:chiral-Gaussian-current-limit}
 \lim_{\delta\downarrow0}
 \E\prod_{i=1}^k
       \delta^{-1}D_\delta^{\varepsilon_i}h_\delta(z_{i,\delta})
   =\sum_{\pi}\ \prod_{\{i,j\}\in\pi}
                  C_{\varepsilon_i,\varepsilon_j}(z_i,z_j),
\end{equation}
where the sum is over pairings, is zero for odd $k$, and
\begin{equation}\label{eq:chiral-contractions}
 C_{++}(z,w)=-\frac{\kappa}{(z-w)^2},\qquad
 C_{--}(z,w)=-\frac{\kappa}{(\bar z-\bar w)^2},\qquad
 C_{+-}(z,w)=C_{-+}(z,w)=0.
\end{equation}
Convergence is uniform on compact subsets of either separated-time
region. By linearity the corresponding assertion holds for all products
of the real centered derivatives $\delta^{-1}D_{x,\delta}$ and
$\delta^{-1}D_{y,\delta}$. These are the off-diagonal derivatives of the
Gaussian covariance
\begin{equation}\label{eq:chiral-log-covariance}
                              K_\kappa(z,w)=-2\kappa\log|z-w|.
\end{equation}
\end{proposition}

\begin{proof}
The preceding lemmas already give the limits as products of pure chiral
functions and their conjugates. It remains to determine $F_k$.
Fix generic spectators $z_2,\ldots,z_k$ as in
Lemma~\ref{lem:chiral-holomorphy}. Its quadratic bound makes every
grid puncture of $F_k$ a pole of order at most two: the Cauchy bound
for a Laurent coefficient of order $-m$, $m>2$, is $Cr^{m-2}$ on a
circle of radius $r$, and tends to zero.

First consider an artificial puncture
$p=p_{jh}=\Re z_j+i\Im z_h$ with $j\ne h$.
Take a fixed small circle $\Gamma$ about its position for the chosen
spectators. In a sufficiently small complex neighborhood of the
spectators, the circle encloses just this moving grid point and remains
in $U_x\cup U_y$. The contour moments
\[
 m_\ell=\frac1{2\pi i}\int_\Gamma
                  \zeta^\ell F_k(\zeta,z_2,\ldots,z_k)\,d\zeta
\]
are jointly holomorphic in the spectator variables, by the uniform
bounds and joint holomorphy on the circle. If the principal part is
$A/(\zeta-p)^2+B/(\zeta-p)$, residue calculation gives
\begin{equation}\label{eq:chiral-pole-moments}
 m_\ell=Bp^\ell+A\ell p^{\ell-1},\qquad
 m_0m_2-m_1^2=-A^2,\qquad
 m_0m_3-m_1m_2=-2pA^2.
\end{equation}
If $A$ is nonzero at any parameter value, the nonzero denominator
persists locally and
\[
             p=\frac{m_0m_3-m_1m_2}
                       {2(m_0m_2-m_1^2)}
\]
is holomorphic there. This is impossible, since
\[
 p=\frac{z_j+\bar z_j}{2}+\frac{z_h-\bar z_h}{2},
                         \qquad \partial_{\bar z_j}p=\tfrac12.
\]
Thus $A$ vanishes. If $B$ is nonzero, $p=m_1/m_0$ gives the same
contradiction. Every artificial puncture is removable.
This removal is joint in the parameters: inside $\Gamma$, the Cauchy
integral of $F_k(\eta,\mathbf z)/(\eta-\zeta)$ supplies a jointly
holomorphic extension and agrees with each one-variable removable
extension.

At a true collision, put
$c=(z_1+z_j)/2$ and $d=z_1-z_j$, retaining generic exterior points.
Locally the union $U_x\cup U_y$ contains every $d\ne0$.
The function
$H(c,d,\mathbf z)=d^2F_k$ is uniformly bounded there and therefore
extends jointly holomorphically to $d=0$. One can see this directly
by its Cauchy coefficients in $d$: boundedness removes the negative
coefficients, uniformly on smaller spectator neighborhoods.
Symmetry under exchanging the two colliding insertions makes $H$ even
in $d$. Lemma~\ref{lem:collision} identifies its constant term as
\[
 H(c,0,\mathbf z)
             =-\kappa F_{k-2}(\mathbf z),
\]
independent of $c$. Subtracting this term leaves an even holomorphic
function with zero constant term, hence a function divisible by $d^2$.
It follows that
\[
 F_k+\frac{\kappa F_{k-2}(\mathbf z)}{(z_1-z_j)^2}
\]
is holomorphic at this diagonal. In particular there is no simple
residue. Independence of the collision center is needed here: a varying
coefficient of $d^{-2}$ could otherwise create a simple term when
$z_j$ is held fixed.

Subtract all these true double-pole terms as a function of $z_1$.
The artificial points are removable, the true points have no remaining
principal parts, and the resulting function is entire.
Lemma~\ref{lem:chiral-holomorphy} gives decay at infinity, and the
subtracted rational terms also decay there. Liouville's theorem yields
\begin{equation}\label{eq:chiral-Wick-recursion}
 F_k(z_1,\ldots,z_k)
   =-\kappa\sum_{j=2}^k
       \frac{F_{k-2}(z_2,\ldots,\widehat z_j,\ldots,z_k)}
                           {(z_1-z_j)^2}.
\end{equation}
Together with $F_0=1$ and $F_1=0$, this is Wick's recursion with
contraction $-\kappa/(z_i-z_j)^2$. Continuity extends it from the
generic spectators to every strict axial chamber. Factorization of
the two chiral sectors and reality then give
\eqref{eq:chiral-Gaussian-current-limit}.

Finally,
$\partial_z\partial_w K_\kappa(z,w)=-\kappa/(z-w)^2$ and
$\partial_z\partial_{\bar w}K_\kappa(z,w)=0$ away from collisions.
The conjugate identity gives the other pure contraction. This proves
the covariance assertion and, by linearity, all the real derivative
formulas. The uniform moving-face convergence established above gives
the asserted local uniformity. Estimates across a time-ordering wall
and for undifferentiated fields will be supplied in
Section~\ref{sec:probability}; none is needed for this proposition.
\end{proof}
